\subsection{Entrenamiento mediante ventana móvil}

Para cada predicción se utilizan los 24 ejemplos supervisados inmediatamente anteriores. Es necesario distinguir la longitud de la historia, de 24 registros por ejemplo, del tamaño del conjunto de entrenamiento, de 24 ejemplos por ajuste. La matriz de entrenamiento tiene dimensión $24\times24$ y el vector de respuestas contiene 24 generaciones.

En el primer ajuste, las respuestas de entrenamiento corresponden a los registros 25 a 48. Cada una se asocia con sus 24 generaciones previas. A continuación, las generaciones de los registros 25 a 48 se utilizan para predecir el registro 49. La primera evaluación requiere, por tanto, 48 observaciones previas. En el ajuste siguiente, las respuestas de entrenamiento abarcan los registros 26 a 49 y se predice el registro 50.

Antes de cada ajuste se estandariza cada columna de entrada mediante \texttt{StandardScaler}, utilizando exclusivamente la media y la desviación estándar del conjunto de entrenamiento. La misma transformación se aplica a la entrada de prueba. La variable de respuesta mantiene su escala original en MWh.

El modelo se ajusta con \texttt{LinearRegression}. Para el instante $t$, se estiman un intercepto y 24 coeficientes mediante
\begin{equation}
(\hat{\beta}_{0,t},\hat{\boldsymbol{\beta}}_t)
=\arg\min_{\beta_0,\boldsymbol{\beta}}
\sum_{s=t-24}^{t-1}
\left(g_s-\beta_0-\widetilde{\boldsymbol{x}}_s^{\mathsf T}
\boldsymbol{\beta}\right)^2,
\end{equation}
donde $\widetilde{\boldsymbol{x}}_s$ representa la entrada estandarizada. La predicción se obtiene como
\begin{equation}
\hat g_t=\hat{\beta}_{0,t}
+\widetilde{\boldsymbol{x}}_t^{\mathsf T}\hat{\boldsymbol{\beta}}_t.
\end{equation}

Después de cada predicción se avanza un registro y se repiten la estandarización y el ajuste. El historial se actualiza con las generaciones observadas disponibles en el nuevo instante. Este procedimiento evalúa pronósticos de un paso y permite adaptar los coeficientes a la ventana más reciente.

Si el subconjunto seleccionado contiene $n$ registros, se construyen $n-24$ ejemplos supervisados y se evalúan $n-48$ predicciones. El programa requiere al menos 50 registros para calcular las métricas sobre dos o más predicciones. Debido a que cada ajuste utiliza 24 ejemplos para estimar 24 coeficientes más un intercepto, la solución puede no ser única. Esta condición debe considerarse al interpretar la estabilidad del modelo.
